\scanpage{23}
\mspage{23}
multiplicatif \deleted{\unsure{autre}} sauf le groupe unité. Comme un sous-groupe d'un groupe de type multiplicatif est de type multiplicatif (Exp VI\quad), il s'ensuit que \deleted{tout} pour tout sous-groupe de type multiplicatif $H$ de $G/Z_0$, on a $H\cap G/Z_0=0$, donc $H\longrightarrow G/G_0$ est injectif. Comme $G/G_0$ est un groupe algébrique radiciel, donc fini sur $k$, cela implique que $H$ est lui-même radiciel et de \uline{rang majoré} par celui de $G/G_0$. Cela \inserted{implique} que la famille des sous-groupes de type multiplicatifs de $G$ admet un plus grand élément, soit $Z$.

Je dis que $G/Z$ n'a pas de sous-groupe de type multiplicatif autre que le sous-groupe unité. Cela résulte du fait (démontré dans Exp VII\quad) qu'une extension commutative de deux groupes algébriques de type multiplicatif est de type multiplicatif; \deleted{compte tenu du caractè-} si donc $Z'$ est l'image inverse dans $G$ d'un sous-groupe de type multiplicatif de $G/Z$, alors $Z'$ est de type multiplicatif, donc $Z'=\unsure{Zc|Z}$ par le caractère maximal de $Z$, donc $Z'/Z=0$.

Il résulte maintenant aisément du ``principe de l'extension finie'' que pour toute extension $K$ de $k$, $(G/Z)_K=G_K/Z_K$ n'a pas non plus de sous-groupe de type multiplicatif sauf le groupe unité.

Nous pouvons maintenant prouver que $Z$ est un centre réductif de $G$. En effet, soit $u:H\longrightarrow G_S$ un homomorphisme de $S$-groupes, où $S$ est un préschéma sur $k$ et $H$ un groupe de type multiplicatif et de type fini sur $S$, prouvons que $u$ se factorise à travers $Z_S$. Il revient au même de dire que l'homomorphisme composé $H\longrightarrow G_S\longrightarrow(G/Z)_S=G_S/Z_S$ est nul. Or je dis que, posant $U=G/Z$, tout homomorphisme $u:H\longrightarrow U_S$ est nul. En effet, en vertu de Exp IX 5.2. on est ramené au cas où $S$ est le spectre d'un corps, et par Exp IX 6.8. cela résulte alors du

\mspage{22}
